\section{Conclusions}
\label{sec:conclusions}

A reconciliation protocol that fixes a rule per class of reference, reports the gap in a unit the reader can use, tests the explanation of a gap and lists what it does not cover, applied to the validation engine of a bank's risk-model validation, gives the following. The engine agrees with \nPrintedAgree\ of \nPrintedN\ printed numbers to the digit and with \nAgreeTwoUnits\ within two units of the last digit; \nBinBeyond\ differ by more than two units, the largest by \nTableNineMaxUnits\ units in Table 9 of Tasche. Three of the six claims of the design document are not met as stated: the count of printed numbers, the closed forms on the extended grid (\nClosedWithin\ of \nClosedCases), and one printed constant of Hull and White (\nWWAtwoUnits\ units from its digit). Of the \nMissTotal\ gaps, \nMissTested\ have an explanation that the experiments test (\nMissRound\ a rounding of the paper's own inputs and \nMissLevel\ the confidence level of the print); the engine agrees with independent recomputations of its tail probabilities to $\nGridJeffMax$ and with QuantLib's closed forms on \nClosedWithin\ of \nClosedCases\ cases at a relative $10^{-6}$; its cohort estimators equal those of \texttt{transitionMatrix} to $\nTMMaxDiff$ with two differences that the package's own code explains; and the exposure of the single swap and of the four-swap netting set agrees with ORE within \nSingleEquivalence\,\% and \nNettingEquivalence\,\% of their effective exposure at the $5\,\%$ level, while the collateralised variants differ from ORE by a gap that repeats on new seeds and that a zero minimum transfer amount only partly explains.

The acceptance rule used for the Monte Carlo reference has a size of \nSizeProdNetting\,\% and a power of \nPowerProdSigmaTwo\,\% against a volatility error of $2\,\%$. Its standard errors overstate the sampling error of antithetic pairs by a factor that grows to \nOverstateLast\ at the later dates, and its thirty-five comparisons would reject two runs of the same model in \nSizePairNetting\,\% of the repetitions with correct errors; the two errors partly cancel on this book and need not elsewhere. Pair-aware errors with the family-wise bound give a size of \nSizeRefineNetting\,\% and a power of \nPowerRefineSigmaTwo\,\%, and a maximum-T bound on pair-aware errors, with a covariance estimated from pilot runs of the engine against itself, gives \nSizePairMaxtNetting\,\% and \nPowerPairMaxtSigmaTwo\,\%, \nMaxtGainNetting\ points above the refinement on the same pairs of runs.

Three steps follow. The tally of printed numbers should be fixed in the product's design document to the one the notebook computes, or the document should state which tables it leaves out. The reconciliation code should adopt the pair-aware error and the family-wise bound, and report the equivalence margin and the detectable gap beside every pass. And the part of the collateralised gap that remains after the zero minimum transfer amount, and the difference of the two dynamic initial margins, should be localised by a run in which ORE's remargining schedule is aligned to the engine's. The independent-route comparison should become the default way to test a tail probability: a test that reuses the engine's helper cannot fail, and the measurements show what the other kind costs.
